\documentclass[12pt]{article}
\usepackage[utf8]{inputenc}
\usepackage[T1]{fontenc}
\usepackage{amsmath,amssymb,amscd}
\usepackage{geometry}
\geometry{margin=2.5cm}

\begin{document}

\begin{center}
\fbox{\textbf{LECTURE 6}}\\[2pt]
\underline{(ALGEBRA)}
\end{center}
\begin{flushright}
\textit{8-XI-94}
\end{flushright}

\bigskip
\noindent\underline{\textbf{The order of an element}}

\bigskip
\noindent\underline{Appendix I}: -- Elements of arithmetic in $\mathbb{N}$.

\bigskip
\noindent\underline{Peano system}: $(A,a,s)$ ; $a\in A$; $s:A\to A$,\footnote{the original reads ``$s:A\to B$''; corrected here to $A\to A$, as required by the definition of a Peano system.} $A\neq\varnothing$

\begin{enumerate}
\item[1)] $s(x)\neq a$ $(\forall)\,x\in A$, \ $\big(s(A)=A\setminus\{a\}\big)$.
\item[2)] $s$ is injective.
\item[3)] If $M\subseteq A$ has the property $\begin{cases} a\in M \\ x\in M \Rightarrow s(x)\in M \end{cases} \Rightarrow M=A$.
\end{enumerate}

\bigskip
\noindent\fbox{T1}. \ If $(A,a,s)$, $(B,b,t)$ are Peano systems $\Rightarrow$

\noindent $\exists!\, f:A\to B$ such that $f(a)=b$, $f\circ s = t\circ f$.

\noindent Moreover,\footnote{the exact phrase is unclear in the original.} $f$ is bijective.

\[
\begin{CD}
A @>s>> A\\
@VfVV @VVfV\\
B @>t>> B
\end{CD}
\]

\bigskip
\noindent Let $(\mathbb{N},0,s)$; \ $s:\mathbb{N}\to\mathbb{N}$.

\noindent the successor function.

\noindent $s(x)=x'$ -- the successor of $x$.

\bigskip
\noindent N.\ Becheanu -- Algebra for Teacher Training\footnote{bibliographic reference; the exact title is reconstructed from a compressed form.}

\bigskip
\noindent\fbox{T2}: $\exists!\, e:\mathbb{N}\times\mathbb{N}\to\mathbb{N}$ \ denoted $e=+$

\noindent $A_1$: $e(m,0)=m$ $(\forall)\,m\in\mathbb{N}$.

\noindent $A_2$: $e(m,n')=e(m,n)'$ $(\forall)\,m,n\in\mathbb{N}$ \quad $m+n'=(m+n)'$.

\bigskip
\noindent It follows that $m+p=m+n \Rightarrow p=n$.

\[
(\mathbb{N},+,0) = \text{a commutative monoid with cancellation.}
\]

\bigskip
\noindent\fbox{T}: $\exists!\, \psi:\mathbb{N}\times\mathbb{N}\to\mathbb{N}$ such that

\noindent $\psi(m,0)=0$\footnote{the original appears to read ``$=m$''; corrected to $0$, as required by the standard definition of multiplication (the adjacent shorthand ``$m\cdot 0$'' is likewise $0$, not $m$).} $(\forall)\,m\in\mathbb{N}$: \quad $m\cdot 0=0$

\noindent $\psi(m,n')=\psi(m,n)+m$ $(\forall)\,m,n\in\mathbb{N}$ \quad $mn'=mn+m$.

\bigskip
\noindent $(\mathbb{N},\cdot,1)$ is a commutative monoid with cancellation.

\noindent $mn=mp \Rightarrow n=p$.

\noindent $m(n+p)=mn+mp$.

\bigskip
\noindent $(\mathbb{N},+,\cdot)$ is an algebraic structure.

\noindent $(\mathbb{N},\leq)$ is an order structure.

\end{document}
